فرض کنید $x[n]$ با دوره‌ی $N$ باشد و نمایش سری فوریه‌ی گسسته‌ی آن به صورت
\[
x[n]=\sum_{k=0}^{N-1} a_k e^{j\frac{2\pi}{N}kn}
\]
باشد. چون $N$ بر $m$ بخش‌پذیر است، می‌نویسیم
\[
N=mL \qquad (L=\tfrac{N}{m}).
\]
طبق شرط مسئله برای هر $n$ داریم:
\begin{flalign*}
0
&=\sum_{r=0}^{L-1} x[n+rL] \\
&=\sum_{r=0}^{L-1}\sum_{k=0}^{N-1} a_k e^{j\frac{2\pi}{N}k(n+rL)} \\
&=\sum_{k=0}^{N-1} a_k e^{j\frac{2\pi}{N}kn}\sum_{r=0}^{L-1} e^{j\frac{2\pi}{N}krL} \\
&=\sum_{k=0}^{N-1} a_k e^{j\frac{2\pi}{N}kn}\underbrace{\sum_{r=0}^{L-1} e^{j\frac{2\pi}{m}kr}}_{C_k}.
&&
\end{flalign*}
حال دو طرف را در $e^{-j\frac{2\pi}{N}\ell n}$ ضرب کرده و روی یک دوره جمع می‌زنیم:
\begin{flalign*}
0
&=\sum_{n=0}^{N-1}\sum_{k=0}^{N-1} a_k C_k e^{j\frac{2\pi}{N}(k-\ell)n}
=\sum_{k=0}^{N-1} a_k C_k \, N\,\delta[k-\ell]
= N a_\ell C_\ell,
&&
\end{flalign*}
پس برای هر $\ell$ داریم $a_\ell C_\ell=0$. اگر $\ell$ بر $m$ بخش‌پذیر باشد، آنگاه
\[
e^{j\frac{2\pi}{m}\ell}=1 \;\;\Longrightarrow\;\;
C_\ell=\sum_{r=0}^{L-1} 1 = L=\frac{N}{m}\neq 0
\]
و در نتیجه
\[
a_\ell=0.
\]
پس برای هر $k$ که بر $m$ بخش‌پذیر باشد، $a_k=0$.